\documentclass[11pt]{article}
\usepackage[T1]{fontenc}
\usepackage[margin=1in]{geometry}
\usepackage{amsmath,amssymb,amsthm}
\usepackage{graphicx}
\usepackage{booktabs}
\usepackage[expansion=false]{microtype}
\usepackage[hidelinks]{hyperref}

\newtheorem{theorem}{Theorem}
\newtheorem{proposition}{Proposition}
\newtheorem{corollary}{Corollary}
\theoremstyle{definition}
\newtheorem{definition}{Definition}
\newtheorem{remark}{Remark}

\newcommand{\dead}{\mathrm{dead}}
\newcommand{\twin}{\mathrm{twin}}

\title{\bf Auditing the Wind Tunnel:\\
the Twin--Enrichment Closed Form at Reachable $\omega$,\\
and the Marginal--versus--Per-Centre Distinction}
\author{Ruqing Chen\\
\small GUT Geoservice Inc., Montreal, Quebec, Canada\\
\small \texttt{ruqing@hotmail.com}}
\date{June 8, 2026\\[4pt]
\small Code, data, and figures: \href{https://github.com/Ruqing1963/6N-windtunnel-audit}{\texttt{github.com/Ruqing1963/6N-windtunnel-audit}}}

\begin{document}
\maketitle

\begin{abstract}
Part XXIII extrapolated the twin--centre enrichment multiplier $\prod_{q\mid N} q/(q-2)$
to the primorial depth $\omega=15$, obtaining $\times 7.03$, and flagged the step honestly as a
closed-form self-portrait assuming local-factor multiplicativity persists past the $\omega\!\approx\!5\text{--}6$
range where it had been checked. We audit that assumption directly. Sieving all centres to
$6N\le 2\times 10^{8}$ ($813{,}370$ twin centres) we measure the actual enrichment, factor by factor.
Three findings. (i) The microscopic \emph{marginal} law is exact: the twin rate among centres divisible
by $q$ exceeds the global rate by $q/(q-2)$, verified to $\lesssim 0.3\%$ for every $q=5,\dots,59$; the
divisor/non-divisor ratio is the complementary $(q-1)/(q-3)$, reconciling the two forms used in Part I
and Part XXIII as two views of one mechanism. (ii) Marginal multiplicativity holds within Poisson error
for every prime combination with an adequate sample ($\lesssim 1\%$ through triples), with a faint
systematic lean to the suppression side---the sign of the Part IV second-order residual, but within our statistics.
(iii) The wind-tunnel object $\prod_{q\mid N}q/(q-2)$, however, is a \emph{marginal} factor and is not the
per-centre likelihood. The correct per-centre enrichment of a centre with divisor set $D$ is
$\mathcal{C}\prod_{q\in D}(q-1)/(q-3)$ with a global constant $\mathcal{C}=\prod_{q>3}\frac{q(q-3)}{(q-1)(q-2)}=0.7216$;
this matches the data to $<1\%$ across the whole reachable range, whereas the wind-tunnel object over-predicts
the per-centre rate by up to $\sim\!38\%$ for typical factorisations. For the squarefree primorial centre the
two nearly coincide ($6.98$ vs.\ $7.03$, a $0.66\%$ gap), so the headline $\times 7.03$ is sound \emph{for the
primorial centre and for the marginal divisibility statement}---but it is an upper envelope, not the typical
$\omega$-stratified enrichment, which sits far below it ($\langle E\rangle\!\approx\!2.6$ at $\omega=6$ versus the
$4.27$ envelope). The wind-tunnel mechanism is thus vindicated microscopically; the two corrections are
interpretive (marginal vs.\ per-centre) and scope (envelope vs.\ typical). As throughout, this is a measurement,
not a theorem, and no claim is made about the infinitude of any constellation.
\end{abstract}

\section{Introduction}
In Volume I the ``wind tunnel'' (Review \S9.1) and the kinetics of Part XXIII synthesised centres
with prescribed prime factorisation and applied the per-factor twin modifier $q/(q-2)$, taking the
product over the $\omega$ prescribed factors to reach an enrichment of $\times7.03$ at the
$\omega=15$ primorial ($N>4.7\times10^{18}$) \cite{PXXIII,review}. Both papers labelled the figure
plainly as a closed-form extrapolation rather than a measurement, and named the load-bearing
assumption: that the local-factor multiplicativity, validated only to $\omega\approx5\text{--}6$,
continues to hold. That is exactly the kind of past extrapolation this programme's honest ledger
asks us to return and test against data. This paper is that audit.

We do not reach $\omega=15$---no sieve does. We reach the $\omega$ that direct enumeration to
$6N\le 2\times10^{8}$ makes statistically solid ($\omega\lesssim 6$), and we test the
\emph{ingredients} of the extrapolation there: the per-prime modifier, its multiplicativity, and---the
point that turns out to matter most---the precise sense in which $\prod q/(q-2)$ is and is not the
enrichment of a given centre.

\section{Method and two notions of enrichment}
For a centre $N$ let $\omega_{>3}(N)$ count the distinct prime factors $>3$ of $N$, and call $N$ a
\emph{twin centre} if $6N-1$ and $6N+1$ are both prime. We sieve all centres $N\le 3.33\times10^{7}$
($6N\le 2\times10^{8}$), computing $\omega_{>3}$ by a distinct-prime sieve and twin membership from a
prime sieve to $2\times10^{8}$; there are $813{,}370$ twin centres, a global rate
$p_0=0.0244011$.

For a prime $q>3$, a twin centre is forbidden on the two killer residues
$\dead(q)=\{6^{-1},-6^{-1}\}\pmod q$; by CRT the survival events at distinct primes are independent,
and the all-centre survival factor at $q$ is $(q-2)/q$. Two distinct conditioning questions then arise,
and conflating them is the whole subject of this paper.

\begin{definition}[Marginal enrichment]
The ratio of the twin rate among centres with $q\mid N$ (all other primes unconstrained) to the global
rate $p_0$. Conditioning on a divisor \emph{set} $D$ marginalises every prime outside $D$.
\end{definition}

\begin{definition}[Per-centre enrichment]
The ratio $\mathbb{E}[\,\text{twin}\mid F\,]/p_0$ where $F$ is the \emph{complete} factorisation of $N$---so
that for \emph{every} prime it is known whether it divides $N$. Binning by $\omega_{>3}(N)$ or by any
function of $F$ measures this notion.
\end{definition}

\section{The microscopic law is exact}
Conditioning on $q\mid N$ fixes $N\equiv0$, a safe residue, so survival at $q$ becomes certain; relative
to the baseline $(q-2)/q$ the factor is $q/(q-2)$. Conditioning instead on $q\nmid N$ leaves $N$ uniform on
the $q-1$ nonzero residues, of which $2$ are killers, giving survival $(q-3)/(q-1)$ and a divisor/non-divisor
ratio $(q-1)/(q-3)$. Table~\ref{tab:perprime} confirms both to $\lesssim0.3\%$ across $q=5,\dots,59$
(Fig.~\ref{fig:audit}A). The $q/(q-2)$ form is the wind-tunnel modifier of Part XXIII; the $(q-1)/(q-3)$ form
is the conditional enrichment tabulated in Part I \cite{PI,PXXIII}. They are not in tension: they are the
baseline-relative and the divisor-relative views of the same two-killer geometry, related identically by
$\frac{q/(q-2)}{(q-1)/(q-3)}=g(q)$ with $g$ as in \S5.

\begin{table}[t]
\centering
\small
\begin{tabular}{rcccc}
\toprule
$q$ & meas.\ $q/(q-2)$ & pred.\ $q/(q-2)$ & meas.\ $(q-1)/(q-3)$ & pred.\ $(q-1)/(q-3)$\\
\midrule
5 & 1.66427 & 1.66667 & 1.99568 & 2.00000 \\
7 & 1.39738 & 1.40000 & 1.49650 & 1.50000 \\
11 & 1.22072 & 1.22222 & 1.24827 & 1.25000 \\
13 & 1.18470 & 1.18182 & 1.20322 & 1.20000 \\
17 & 1.13888 & 1.13333 & 1.14885 & 1.14286 \\
19 & 1.11622 & 1.11765 & 1.12347 & 1.12500 \\
23 & 1.09346 & 1.09524 & 1.09812 & 1.10000 \\
29 & 1.07675 & 1.07407 & 1.07971 & 1.07692 \\
37 & 1.05909 & 1.05714 & 1.06083 & 1.05882 \\
47 & 1.04376 & 1.04444 & 1.04475 & 1.04545 \\
59 & 1.04019 & 1.03509 & 1.04091 & 1.03571 \\
\bottomrule
\end{tabular}
\caption{Per-prime marginal enrichment at $6N\le2\times10^{8}$. Both closed forms hold to $\lesssim0.3\%$;
deviations change sign with $q$ and are consistent with statistical noise.}
\label{tab:perprime}
\end{table}

\section{Multiplicativity holds within statistics}
The extrapolation's core assumption is that distinct primes' modifiers multiply. Under CRT independence this
is exact for the marginal enrichment: conditioning on a divisor set $D$ gives
$\prod_{q\in D}q/(q-2)$. Table~\ref{tab:mult} tests it. Every combination with an adequate sample agrees with
the product to $\lesssim1\%$; the apparent $-8.5\%$ at $\{5,7,11,13\}$ rests on only $\sim\!500$ twins
(relative Poisson error $\approx4.5\%$, so $\sim\!2\sigma$). The residuals carry a faint systematic lean to
the suppression side, the sign of the second-order inter-centre residual isolated empirically in Part IV
\cite{PIV}, but they do not exceed $\sim\!2\sigma$ anywhere; within the reachable range, multiplicativity is
confirmed and the wind-tunnel product is the right marginal object.

\begin{table}[t]
\centering
\small
\begin{tabular}{lrccc}
\toprule
divisor set $D$ & $n$ centres & meas.\ enrich & $\prod_{q\in D}\frac{q}{q-2}$ & resid.\ (\%)\\
\midrule
$\{5,7\}$ & 952380 & 2.3256 & 2.3333 & $-0.33$ \\
$\{5,11\}$ & 606060 & 2.0272 & 2.0370 & $-0.48$ \\
$\{7,11\}$ & 432900 & 1.6944 & 1.7111 & $-0.98$ \\
$\{5,13\}$ & 512820 & 1.9712 & 1.9697 & $+0.07$ \\
$\{5,7,11\}$ & 86580 & 2.8424 & 2.8519 & $-0.33$ \\
$\{5,7,13\}$ & 73260 & 2.7109 & 2.7576 & $-1.69$ \\
$\{5,7,11,13\}$ & 6660 & 3.0829 & 3.3704 & $-8.53$ \\
\bottomrule
\end{tabular}
\caption{Marginal multiplicativity. Agreement is $\lesssim1\%$ wherever the twin sample exceeds
$\sim\!10^{4}$; the last row is sample-starved ($\sim\!500$ twins).}
\label{tab:mult}
\end{table}

\section{The marginal--per-centre distinction (the substantive correction)}
The wind tunnel reports $\prod_{q\mid N}q/(q-2)$ as ``the enrichment of a centre'' with that factorisation.
But $q/(q-2)$ is a \emph{marginal} factor, and a per-centre likelihood must also account for the primes that
do \emph{not} divide $N$: knowing the complete factorisation $F$ says that each $q\nmid N$ sits at a uniformly
random \emph{nonzero} residue, safe with probability $(q-3)/(q-1)$, not $(q-2)/q$. Writing $D$ for the divisor
set,
\[
\frac{\mathbb{E}[\twin\mid F]}{p_0}
=\prod_{q\in D}\underbrace{\frac{1}{(q-2)/q}}_{q/(q-2)}\ \cdot\
\prod_{q\notin D}\underbrace{\frac{(q-3)/(q-1)}{(q-2)/q}}_{g(q)},
\qquad
g(q):=\frac{q(q-3)}{(q-1)(q-2)}<1 .
\]
Let $\mathcal{C}:=\prod_{q>3}g(q)$. Since $\prod_{q\notin D}g(q)=\mathcal{C}/\prod_{q\in D}g(q)$ and
$\frac{q/(q-2)}{g(q)}=\frac{q-1}{q-3}$, the divergent-looking expression collapses to a finite per-centre law:
\begin{equation}
\boxed{\ \frac{\mathbb{E}[\twin\mid F]}{p_0}=\mathcal{C}\prod_{q\mid N}\frac{q-1}{q-3}\ },
\qquad
\mathcal{C}=\prod_{q>3}\frac{q(q-3)}{(q-1)(q-2)}=0.72160 .
\label{eq:percentre}
\end{equation}
Each factor of $\mathcal{C}$ is $1-\tfrac{2}{(q-1)(q-2)}=1-O(1/q^2)$, so the product converges absolutely---the
same $O(1/q^2)$ floor that makes the lattice-correlation products of Part XXII converge. The constant $\mathcal{C}$
is a genuine global suppression carried by the primes that miss $N$, and it is what the bare product
$\prod q/(q-2)$ omits.

Equation~\eqref{eq:percentre} is not a fit: it is an exact CRT consequence, and it reproduces the data
(Fig.~\ref{fig:audit}B). Binning centres by their per-centre model value
$\mathcal{C}\prod_{q\mid N}(q-1)/(q-3)$, the measured per-centre enrichment lies on the diagonal to within
$1\%$ across the entire range, while the wind-tunnel object $\prod q/(q-2)$ falls systematically below the
diagonal---over-predicting the per-centre rate by a factor $\mathcal{C}\prod_{q\mid N}g(q)^{-1}$ that runs
from $\mathcal{C}=0.72$ for factor-poor centres up to $0.94$ for small-prime-rich ones, exactly as
\eqref{eq:percentre} requires.

\begin{figure}[t]
\centering
\includegraphics[width=\textwidth]{fig1_audit.pdf}
\caption{Auditing the closed form against real centres. (A) The microscopic marginal law: measured
enrichment versus prediction, $q=5,\dots,59$, lies on $y=x$ for both the baseline-relative $q/(q-2)$ and the
divisor-relative $(q-1)/(q-3)$ forms; residuals (inset) are $\pm0.5\%$. (B) Per-centre enrichment versus model:
the correct law $\mathcal{C}\prod(q-1)/(q-3)$ (green) sits on the diagonal, while the wind-tunnel object
$\prod q/(q-2)$ (orange) systematically over-predicts the per-centre rate.}
\label{fig:audit}
\end{figure}

\section{The $\times7.03$ verdict}
Two facts now sit cleanly side by side.

\paragraph{It is correct for the primorial centre.} For the squarefree primorial $\Omega_{15}$ (exactly the
factors $\{5,7,\dots,59\}$, nothing else), \eqref{eq:percentre} gives
$\mathcal{C}\prod_{q\le59}(q-1)/(q-3)=6.983$, against the wind-tunnel $\prod_{q\le59}q/(q-2)=7.030$: a gap of
only $-0.66\%$ (Fig.~\ref{fig:envelope}B). The two nearly coincide because of an exact identity---the
per-centre value equals $\big[\prod_{q\le59}q/(q-2)\big]\cdot\prod_{q>59}g(q)$, and the tail
$\prod_{q>59}g(q)\approx0.9934$ almost exactly cancels the difference between the $(q-1)/(q-3)$ boost on the
fifteen small primes and the suppression $\mathcal{C}$. For the primorial centre itself, then, $\times7.03$ is
right to better than a percent. It is equally right as the \emph{marginal} statement ``centres divisible by
$\{5,\dots,59\}$ are $7.03\times$ more likely to be twin centres,'' which is the multiplicative product
confirmed in \S4.

\paragraph{But it is an upper envelope, not the typical enrichment.} The primorial centre is vanishingly rare.
A \emph{typical} centre at a given $\omega$ does not carry the smallest primes; it carries a mixture, mostly
larger ones whose modifiers are close to $1$. The measured average enrichment $\langle E\rangle(\omega)$ therefore
sits far below the primorial envelope $\prod_{i\le\omega}q_i/(q_i-2)$ at matched $\omega$
(Fig.~\ref{fig:envelope}A): at $\omega=6$ the data give $2.63$ against the envelope's $4.27$. The correct
model average $\langle\mathcal{C}\prod(q-1)/(q-3)\rangle$ tracks the data through $\omega=5$
(Table~\ref{tab:omega}); the $\omega=6$ bin, with $190$ twins, is sample-starved (the $-20\%$ there is
$\sim\!3\sigma$ and consistent with either noise or a faint extra suppression). The headline $\times7.03$ is
thus the value reached only by the primorial-loaded subset of measure tending to zero, not the enrichment of
the generic $\omega=15$ centre.

\begin{table}[t]
\centering
\small
\begin{tabular}{rrrccc}
\toprule
$\omega$ & $n$ centres & twins & meas.\ $\langle E\rangle$ & model \eqref{eq:percentre} & primorial $\prod\frac{q}{q-2}$\\
\midrule
1 & 6,741,131 & 120,175 & 0.731 & 0.722 & 1.667 \\
2 & 14,241,165 & 311,377 & 0.896 & 0.891 & 2.333 \\
3 & 9,700,927 & 276,430 & 1.168 & 1.171 & 2.852 \\
4 & 2,449,886 & 94,528 & 1.581 & 1.614 & 3.370 \\
5 & 197,049 & 10,654 & 2.216 & 2.297 & 3.820 \\
6 & 2,958 & 190 & 2.632 & 3.313 & 4.269 \\
\bottomrule
\end{tabular}
\caption{$\omega$-stratified enrichment. The measured average tracks the per-centre model
\eqref{eq:percentre}, both far below the primorial envelope at matched $\omega$. The last row is
sample-starved.}
\label{tab:omega}
\end{table}

\begin{figure}[t]
\centering
\includegraphics[width=\textwidth]{fig2_envelope.pdf}
\caption{The $\times7.03$ figure in context. (A) Measured average enrichment $\langle E\rangle(\omega)$ (blue,
with Poisson bars) and the correct per-centre model (green) sit far below the primorial envelope (purple),
which reaches $\times7.03$ only at the $\omega=15$ primorial. (B) For the squarefree primorial centre itself,
the wind-tunnel product ($7.030$) and the correct per-centre value ($6.983$) agree to $0.66\%$.}
\label{fig:envelope}
\end{figure}

\section{Scope and honest ledger}
The audit returns a clear verdict on a past extrapolation. The wind-tunnel \emph{mechanism} is vindicated:
the per-prime modifier $q/(q-2)$ is exact to $\lesssim0.3\%$, and its multiplicativity holds within Poisson
error throughout the reachable range. The headline number $\times7.03$ is sound for what it actually
describes---the squarefree primorial centre (to $0.66\%$) and the marginal divisibility statement (exactly,
under multiplicativity). The audit adds two clarifications and no new theorem. First, an \emph{interpretive}
one: $q/(q-2)$ is a marginal factor; the per-centre likelihood of a centre with a given factorisation is
$\mathcal{C}\prod(q-1)/(q-3)$ with $\mathcal{C}=0.7216$, and using the bare product as a per-centre multiplier
over-predicts non-primorial factorisations by up to $\sim\!38\%$. This is also the clean reconciliation of the
Part I $(q-1)/(q-3)$ form with the Part XXIII $q/(q-2)$ form, now seen as the per-centre and marginal views of
the identical two-killer geometry, joined by $g(q)$ and the constant $\mathcal{C}$. Second, a \emph{scope} one:
$\times7.03$ is an upper envelope, realised by a vanishing subset, not the typical $\omega=15$ enrichment.

We state the limits plainly. The data reach $6N\le2\times10^{8}$, smaller than the $6\times10^{8}$ run of
Part XXVI \cite{PXXVI} but ample for per-prime and per-centre statistics through $\omega\approx6$; the
behaviour at $\omega=15$ remains an extrapolation, now built on ingredients measured to hold exactly where we
can see them. The faint suppression lean in Tables~\ref{tab:mult} and~\ref{tab:omega} is the right sign for
the Part IV second-order residual and the near-Poisson tilt of Part XXVI but does not exceed $\sim\!2\text{--}3\sigma$;
isolating it cleanly is a second-moment question we leave where Part XXVI left it. Everything here equals the
discrete sieve count, takes the Hardy--Littlewood heuristic as input, and makes no claim about the infinitude
of twins or of any constellation. It is an audit of our own self-portrait, and the portrait holds.

\paragraph{Data and code availability.} The sieve experiment (to $6N=2\times10^{8}$), the per-prime,
multiplicativity, per-centre and $\omega$-stratified computations, the two figures, and a verifier of the
per-centre identity \eqref{eq:percentre} and the constant $\mathcal{C}$ are openly available at
\href{https://github.com/Ruqing1963/6N-windtunnel-audit}{\texttt{github.com/Ruqing1963/6N-windtunnel-audit}}.

\begin{thebibliography}{9}
\bibitem{HL} G.~H.~Hardy, J.~E.~Littlewood, \emph{Some problems of `Partitio Numerorum'; III}, Acta Math.\ \textbf{44} (1923), 1--70.
\bibitem{EK} P.~Erd\H{o}s, M.~Kac, \emph{The Gaussian law of errors in the theory of additive number theoretic functions}, Amer.\ J.\ Math.\ \textbf{62} (1940), 738--742.
\bibitem{PI} R.~Chen, \emph{Conditional density on the 6N skeleton} (Part I), Zenodo, doi:10.5281/zenodo.20470367 (2026).
\bibitem{PIV} R.~Chen, \emph{Second-order spatial-compression residual} (Part IV), Zenodo, doi:10.5281/zenodo.20500465 (2026).
\bibitem{PXII} R.~Chen, \emph{The microscopic root $1/(q-2)$ and the mean shift} (Part XII), Zenodo, doi:10.5281/zenodo.20528446 (2026).
\bibitem{PXVII} R.~Chen, \emph{Arithmetic geodynamics on the 6N skeleton} (Part XVII), Zenodo, doi:10.5281/zenodo.20541575 (2026).
\bibitem{PXXIII} R.~Chen, \emph{Phase transition kinetics and topological event horizons at extreme $\omega$} (Part XXIII), Zenodo, doi:10.5281/zenodo.20586919 (2026).
\bibitem{PXXVI} R.~Chen, \emph{The second moment of $\omega$ under twin-centre conditioning on the 6N skeleton} (Part XXVI), Zenodo, doi:10.5281/zenodo.20593654 (2026).
\bibitem{review} R.~Chen, \emph{Arithmetic geodynamics on the 6N skeleton: a review of Parts I--XIX}, Zenodo, doi:10.5281/zenodo.20585301 (2026).
\end{thebibliography}

\end{document}
